\documentclass[]{article}

\usepackage{amsmath}
\usepackage{multirow}
\usepackage{natbib}
\usepackage{graphicx} 


\begin{document}

In this paper, we investigated the anomalous transport of passive tracers in a three-dimensional, quenched velocity field generated by a static configuration of vortex filaments. The primary goal was to connect the theoretical prediction of superdiffusion governed by Lévy-stable Holtsmark statistics to the transport dynamics observed in a finite-density system and to elucidate the underlying physical mechanisms. To this end, we performed direct numerical simulations of tracer trajectories across a range of filament densities and analyzed the resulting transport statistics, including the mean squared displacement, velocity distributions, and velocity correlations, alongside the local topology of the flow field.

Our results confirm that the transport is strongly superdiffusive. The anomalous diffusion exponent, $\alpha$, systematically decreases from a nearly ballistic value ($\alpha \approx 2$) at low filament densities towards the theoretically predicted asymptotic value of $\alpha = 1.5$ as the system becomes more crowded. However, this convergence is slow, and even at the highest density studied, the measured exponent ($\alpha \approx 1.84$) remains significantly above the theoretical limit. The underlying velocity statistics were found to be consistent with the Holtsmark distribution, exhibiting the predicted heavy-tailed, power-law decay. This confirms that the velocity field possesses the Lévy-stable character required for this class of anomalous transport.

From these results, we have learned several key aspects of transport in this system. First, a direct mechanistic link exists between the flow's geometric structure and the emergent transport dynamics. We demonstrated that low-speed trapping events, which contribute significantly to the anomalous diffusion, are localized in rotation-dominated regions of the flow, whereas high-speed flights occur in strain-dominated regions. The residence times in these trapping zones follow a heavy-tailed distribution, providing a physical origin for the Lévy-like statistics. Second, we have learned that the transport is fundamentally shaped by persistent velocity correlations inherent to the quenched, deterministic flow field. This memory, evidenced by the long correlation times and anti-correlation signatures in the velocity autocorrelation function, distinguishes the process from a canonical, memoryless Lévy walk. These strong correlations are responsible for the slow convergence of the diffusion exponent to its asymptotic value. Finally, the system exhibits signatures of non-ergodicity, with significant variability between individual tracer trajectories, highlighting that transport in such quenched, disordered media cannot be fully described by ensemble averages alone. In summary, this work provides a comprehensive picture of how quenched spatial disorder in a three-dimensional vortex field generates a distinct, non-ergodic superdiffusive regime, bridging the gap between abstract statistical theory and concrete physical dynamics.

\end{document}
                